%% =====================================================================
%%  Corrigés — Bilan 3 « En ville »
%% =====================================================================
\corrigesBilan{3}

\begin{corrige}{1}
\begin{reponses}[2]
\item \rep{Turn left.}
\item \rep{Go straight ahead.}
\item \rep{Cross the street.}
\item \rep{Don't take the bus.}
\end{reponses}
\piege{« tout droit » se dit \en{straight ahead}, et non \en{right} (« à droite »).}
\end{corrige}

\begin{corrige}{2}
\begin{reponses}[3]
\item \rep{How much}
\item \rep{How many}
\item \rep{some}
\item \rep{any}
\item \rep{any}
\item \rep{How much}
\end{reponses}
\pourquoi{\en{How many} devant un pluriel (\en{apples}) ; \en{How much} pour un prix ou un indénombrable
(\en{money}). \en{some} dans une affirmation, \en{any} dans une négation ou une question.}
\end{corrige}

\begin{corrige}{3}
\begin{reponses}[2]
\item \rep{fifteen dollars}
\item \rep{fifty dollars}
\item \rep{thirteen ninety-nine}
\item \rep{eighty dollars}
\end{reponses}
\retiens{pour un prix avec des cents, on dit souvent les deux nombres sans le mot \en{dollars} :
\en{\$13.99} = \en{thirteen ninety-nine}.}
\end{corrige}

\begin{corrige}{4}
\begin{reponses}[3]
\item \rep{order}
\item \rep{Can}
\item \rep{Would}
\item \rep{like}
\item \rep{are}
\item \rep{check}
\end{reponses}
\pourquoi{\en{Here you are} accompagne un objet que l'on tend. Au Royaume-Uni, l'addition se dit
\en{the bill}.}
\end{corrige}

\begin{corrige}{5}
\begin{reponses}[4]
\item \rep{can}
\item \rep{can't}
\item \rep{can't}
\item \rep{can}
\end{reponses}
\end{corrige}

\begin{corrige}{6}
\begin{reponses}[3][french]
\item \rep{V}
\item \rep{F}
\item \rep{V}
\item \rep{F}
\item \rep{V}
\item \rep{F}
\end{reponses}
\pourquoi{Léa fait du M (\en{medium}) ; le pull coûte 39 dollars (\en{thirty-nine}) ; elle paie par carte.
Texte de la piste page~\pageref{trans-B3.01}.}
\end{corrige}

\begin{corrige}{7}
\begin{reponses}[3][french]
\item \rep{V}
\item \rep{V}
\item \rep{F}
\item \rep{V}
\item \rep{F}
\item \rep{F}
\end{reponses}
\pourquoi{le steak coûte 24 dollars ; il n'y a que deux desserts ; \en{Tax and tip not included} : la
taxe et le pourboire ne sont pas compris.}
\end{corrige}

\begin{corrige}{8}
Réponse modèle :
\begin{anglais}
\begin{tcolorbox}[otBase,colback=otGreenTint,colframe=otGreenTint,fontupper=\fontsize{10.4}{14}\selectfont]
I can swim quite well. I can cook very well. I can play the piano, but not very well. I can't ski at
all!
\end{tcolorbox}
\end{anglais}
\retiens{0,5 point par phrase correcte : \en{can} / \en{can't} + base verbale, adverbe bien placé.}
\end{corrige}

\begin{corrige}{9}
Réponses modèles :
\begin{reponses}[1]
\item \rep{Can I have a latte, please?}
\item \rep{Yes, please. I'd like a muffin.}
\item \rep{To go, please.} \textit{ou} \rep{For here, please.}
\item \rep{Can I pay by card?}
\item \rep{Thank you! You too! Bye!}
\end{reponses}
\end{corrige}
